\section{Conclusiones}

Las conclusiones de este documento recapitulan los aspectos centrales descritos en las secciones anteriores, evaluando el alcance funcional del Producto Mínimo Viable (MVP) entregado, el ecosistema documental construido y el impacto real de la plataforma. Esta visión holística permite dimensionar el verdadero profesionalismo del proyecto y cómo este materializa la formación integral recibida a lo largo de la carrera.

\subsection{Resultados principales: Ecosistema organizativo y repositorios}

Absolutamente toda la implementación técnica, lógica y documental del Trabajo Profesional reside de forma centralizada en nuestra organización oficial de GitHub: \url{https://github.com/TrabajoProfesional-AGGZ}. La magnitud del proyecto exigió una arquitectura sumamente distribuida, componiéndose de \textbf{27 repositorios y 2 proyectos Kanban}.

\subsubsection{Gestión ágil y control de incidencias (Proyectos Kanban)}
\begin{itemize}
    \item \textbf{Centro de monitoreo de tareas:} Tablero estructurado bajo la metodología Kanban que rigió la ejecución de los 13 \textit{sprints} de desarrollo. Permitió gestionar el \textit{backlog} de manera dinámica, asignando responsables, pesos (esfuerzo) y prioridades a cada historia de usuario.
    \item \textbf{Centro de monitoreo de bugs:} Espacio homólogo dedicado exclusivamente a la detección, seguimiento y resolución de errores (\textit{bugs}) y deudas técnicas reportadas por los auditores automáticos (SonarCloud/OpenClaw).
\end{itemize}

\subsubsection{Repositorios privados (Lógica de negocio e infraestructura)}
Con el fin de resguardar la propiedad intelectual, la seguridad criptográfica y los derechos de autor de las implementaciones subyacentes, los repositorios \textit{back-end} se mantienen en carácter privado. Sin embargo, \textbf{ante cualquier requerimiento académico, de evaluación o interés técnico justificado, el equipo está totalmente a disposición para otorgar los accesos correspondientes}. El código fuente de estos repositorios está documentado internamente siguiendo los más altos estándares y buenas prácticas (\textit{self-documenting code}):

\begin{itemize}
    \item \texttt{microservicio-club}: Gestiona la lógica central (socios, disciplinas, reservas y eventos).
    \item \texttt{microservicio-pagos}: Orquesta las transacciones financieras y pasarelas externas.
    \item \texttt{microservicio-autenticacion}: Administra la identidad, seguridad y roles de los usuarios.
    \item \texttt{microservicio-acceso}: Controla la validación criptográfica (TOTP) y lectura de QRs.
    \item \texttt{microservicio-analiticas}: Calcula y distribuye métricas y proyecciones de morosidad.
    \item \texttt{microservicio-bot-conversacional}: Contiene la lógica asíncrona del bot de Telegram.
    \item \texttt{gateway}: \textit{API Gateway} que unifica y protege el acceso a la red interna.
    \item \texttt{healthchecker}: Monitor de disponibilidad global del ecosistema y de los clubes clientes.
    \item \texttt{openclaw}: Agente de QA especializado para testing intensivo de ciberseguridad.
    \item \texttt{toolkit-financiero}: Herramientas de cálculo, proyección y análisis comercial.
\end{itemize}

\subsubsection{Repositorios públicos (Interfaces y portales documentales)}

La totalidad de las interfaces cliente y la documentación asociada a los repositorios privados se mantienen en carácter público para demostrar la transparencia, alcance y arquitectura del proyecto sin comprometer credenciales sensibles.

\textbf{Portales documentales de la arquitectura (\textit{Back-end}):} Para cada microservicio privado existe un repositorio espejo terminado en \texttt{-doc}. Estos alojan portales generados con JustTheDocs que explican exhaustivamente la arquitectura del servicio original: diagramas de contexto, justificación tecnológica, métricas de implementación y el listado de \textit{endpoints} públicos:
\begin{itemize}
    \item \texttt{microservicio-club-doc}, \texttt{microservicio-pagos-doc}, \texttt{microservicio-autenticacion-doc}, \texttt{microservicio-acceso-doc}, \texttt{microservicio-analiticas-doc}, \\ \texttt{microservicio-bot-conversacional-doc}, \texttt{gateway-doc}, \texttt{healthchecker-doc}.
\end{itemize}

\textbf{Interfaces gráficas (\textit{Front-end}):} Exponen el código y sus propias páginas documentales con vistas principales del sistema.
\begin{itemize}
    \item \texttt{plataforma-web}: Panel de administración (\textit{dashboard}) para los directivos.
    \item \texttt{aplicacion}: Aplicación móvil (PWA) interactiva para los socios.
    \item \texttt{control-de-accesos}: Aplicación móvil (PWA) para el escaneo de QRs por parte de empleados.
\end{itemize}

\textbf{Repositorios documentales / Páginas de gestión global:}
\begin{itemize}
    \item \texttt{.github}: Centro neurálgico que conecta repositorios y presenta al equipo.
    \item \texttt{scrum}: Documentación de las ceremonias, métricas y evolución de los 13 \textit{sprints}.
    \item \texttt{bitacora}: Reportes detallados del avance a lo largo de las 37 semanas del proyecto.
    \item \texttt{ideas-productos}: El proceso creativo y analítico (5Cs, 4Ps) de las 13 ideas primigenias.
    \item \texttt{documentacion}: El repositorio maestro que compila todos los entregables académicos.
\end{itemize}

\subsection{El inmenso valor documental del proyecto}

El hito más destacable de SocioUnido, a la par de su alcance tecnológico, es su aporte documental. Se han construido \textbf{16 páginas plenamente documentales con JustTheDocs} que abarcan el Trabajo Profesional de punta a punta (8 de arquitectura, 3 visuales y 5 de gestión global). Particularmente, el repositorio \texttt{documentacion} aloja una vasta cantidad de material complementario:

\begin{itemize}
    \item El Anteproyecto y las versiones iterativas de la Entrega Intermedia.
    \item Evidencia de validación del producto (entrevistas a trabajadores y encuestas de usabilidad).
    \item El Informe Final definitivo y diagramas del Modelo C4 / DER.
    \item Redacciones de discursos comerciales (\textit{Pitch}) y desarrollo de identidad visual de marca.
    \item Más de \textbf{25 anexos formales} analizando aristas técnicas, legales y de gestión.
    \item Decenas de videotutoriales explicativos de la plataforma con sus manuales de uso.
\end{itemize}

\subsection{Identidad comercial, innovación y contribución a otras entidades}

Desde su concepción, SocioUnido fue tratado como un producto real. Para respaldar esta visión, se consolidó una \textbf{Página publicitaria (Landing Page)} (\url{https://sociounido.com/}) y un \textbf{Canal de YouTube} (\url{https://www.youtube.com/@SocioUnido}) para validación y tutoriales. Este enfoque comercial fue validado de forma empírica y constante mediante reuniones con trabajadores de clubes reales y encuestas a usuarios finales.

A partir de este desarrollo, el Trabajo Profesional ha generado \textbf{innovaciones tecnológicas} significativas que contribuyen directamente a la modernización, inclusión y saneamiento de las instituciones deportivas (especialmente en las categorías de ascenso):
\begin{enumerate}
    \item \textbf{Innovación criptográfica (Infraestructura \textit{offline}):} Adaptación del algoritmo TOTP para la validación de accesos en estadios sin conexión a internet, mitigando el colapso de las redes 4G/5G en eventos masivos y erradicando el fraude sin requerir hardware costoso.
    \item \textbf{Innovación predictiva:} Desarrollo de herramientas para predecir el riesgo de fuga societaria basándose en historiales y métricas definidas por el club, transformando la cobranza reactiva en un modelo de retención proactiva.
    \item \textbf{Innovación en accesibilidad (PLN):} Democratización del sistema mediante un asistente omnicanal en Telegram con Procesamiento de Lenguaje Natural, eliminando la fricción tecnológica para usuarios de todas las edades.
\end{enumerate}

\subsection{Impacto en la formación: Conocimientos aplicados}

El diseño, arquitectura y desarrollo de la plataforma SocioUnido no es un esfuerzo aislado, sino la síntesis y aplicación directa de los conocimientos adquiridos a lo largo de toda la carrera de Ingeniería en Informática. La resolución de las problemáticas de las instituciones deportivas exigió un enfoque multidisciplinario:

\begin{itemize}
    \item \textbf{Fundamentos matemáticos, ciencias básicas y algoritmos:} \textit{Física} proveyó el conocimiento sobre la propagación de señales, justificando la validación de ingresos de forma \textit{offline} en estadios repletos. \textit{Álgebra A y Álgebra Lineal} fueron indispensables para comprender la criptografía (aritmética modular y \textit{hash}) detrás de la innovación TOTP. \textit{Química} aportó la perspectiva de las limitaciones de las baterías de los dispositivos de escaneo móvil a la intemperie. \textit{Algoritmos y Estructuras de Datos, y Teoría de Algoritmos} constituyeron el núcleo de eficiencia, optimizando las búsquedas y el filtrado de enormes padrones. Finalmente, \textit{Análisis Matemático y Probabilidad y Estadística} aportaron el marco teórico para la evaluación y limpieza de datos.
    
    \item \textbf{Ingeniería de software, arquitectura y calidad:} \textit{Ingeniería de Software I y II} guiaron la concepción del ciclo de vida, la aplicación de Scrum, arquitecturas de microservicios (alta cohesión y bajo acoplamiento) y el \textit{Delivery Pipeline} (CI/CD). Materias como \textit{Taller de Programación y Paradigmas de Programación} forjaron las bases de pruebas automatizadas (TDD) y el desarrollo de código limpio y escalable.

    \item \textbf{Sistemas, redes y bases de datos:} \textit{Organización del Computador y Sistemas Operativos} permitieron comprender la arquitectura a bajo nivel y el manejo de memoria en contenedores. \textit{Base de Datos} aseguró la persistencia y normalización mediante transacciones concurrentes en PostgreSQL. \textit{Programación Concurrente, Redes y Sistemas Distribuidos I} resultaron la columna vertebral de la plataforma Cloud, garantizando balanceo de carga, protocolos de red (HTTP/TCP), APIs RESTful y tolerancia a fallos.

    \item \textbf{Ciencia de datos y Ciberseguridad:} Los conocimientos de \textit{Ciencia de Datos} sirvieron para procesar grandes volúmenes de transacciones y extraer características proyectando soluciones de retención. En paralelo, el \textit{Taller de Seguridad Informática} aportó las técnicas de defensa en profundidad, sanitización de datos y gestión de roles (JWT) para proteger la información.

    \item \textbf{Emprendedurismo, gestión y contexto profesional:} \textit{Introducción al Conocimiento de la Sociedad y el Estado e Introducción al Pensamiento Científico} brindaron el pensamiento crítico para entender la dinámica sociopolítica de los clubes. \textit{Empresas de Base Tecnológica I y II} sentaron las bases corporativas, financieras y legales del modelo SaaS B2B, estrategias de \textit{pricing} y propiedad intelectual. Por último, \textit{Gestión del Desarrollo de Sistemas Informáticos} aportó la estructura ejecutiva para estimar tiempos, liderar al equipo y gestionar los riesgos del proyecto.
\end{itemize}

\subsection{Cierre y reflexión final}

En síntesis, SocioUnido no es un mero ensamblaje de tecnologías, sino el producto maduro de un ingeniero informático formado integralmente en la Universidad de Buenos Aires (UBA), capaz de integrar ciencias básicas, aplicadas, y competencias de gestión para resolver problemas complejos de la sociedad y la industria.

El código fuente en su totalidad cuenta con más de \textbf{1.200} \textit{commits} registrados y \textbf{290} revisiones cruzadas (\textit{Pull Requests}) evaluadas. El volumen estructural desarrollado impactó en aproximadamente \textbf{400.000} líneas de código a lo largo de este año de desarrollo continuo.

Por esto, y por el impacto real del producto construido, el equipo no podría estar más orgulloso de los resultados obtenidos. Tenemos la profunda convicción de que el Trabajo Profesional que estamos entregando representa el más alto estándar de excelencia institucional. Haber fortalecido el proyecto en todos los aspectos evaluables (tecnológicos, comerciales, documentales y arquitectónicos) demuestra que el trayecto formativo de nuestra carrera ha rendido sus frutos, cerrando esta etapa universitaria con un broche de oro materializado en este impresionante entregable final.
